% ============================================================
% Feldman 2020 positioning — two versions:
%   (A) for Introduction (shorter)
%   (B) for Related Work (longer, more technical)
% ============================================================

% ── (A) Introduction version ────────────────────────────────
% Insert after stating the research question, before contributions.
%
Feldman~\cite{feldman2020memorization} showed that \emph{individual} rare
samples require memorization to achieve good generalization, and that
memorization is observable only after training.
Our work asks a complementary and temporally prior question: can
\emph{dataset-level} structural statistics predict \emph{aggregate}
membership inference risk before any model is trained?
Where Feldman identifies which samples will be memorized (post-training,
sample-level), we predict how much a dataset will leak (pre-training,
dataset-level)---a distinction that is both scientifically and
practically fundamental.


% ── (B) Related Work version ────────────────────────────────
% Insert as a paragraph in the Related Work section.
%
\paragraph{Memorization and long-tail samples.}
Feldman~\cite{feldman2020memorization} proved that learning algorithms
must memorize atypical (long-tail) training examples to achieve near-optimal
generalization, and connected this to per-sample privacy risk.
Subsequent work~\cite{carlini2022lira,zhang2021counterfactual} confirmed
that samples with high training loss or low density are individually more
vulnerable to membership inference.
Our work differs along three dimensions.
\emph{Unit of analysis:} we study \emph{dataset-level} statistics rather
than per-sample certificates.
\emph{Timing:} we operate \emph{before} training, producing a risk estimate
without a trained model.
\emph{Output:} we produce a deployable scalar index (DPRI) that
practitioners can compute at data-collection time, not a post-hoc
explanation of memorization.
Concretely, our experiments include datasets where \emph{no individual
sample is rare} (e.g., high-density synthetic populations), yet DPRI
correctly predicts elevated aggregate risk---a regime where sample-level
memorization theory offers no prediction.
